\section*{摘 要}
\addcontentsline{toc}{section}{摘 要}

本报告以凯尔纳（Kellner）目镜为研究对象，围绕目视光学系统的小型化、较大视场和像质均衡要求，在Zemax OpticStudio软件平台上完成了目镜的初始结构建立、评价函数设置、自动优化、性能分析与公差评估。凯尔纳目镜是一种经典的三片式目镜结构，通常由一片场镜和一组消色差双合镜构成，具有结构简单、成本较低、轴上色差校正较好的特点，适用于入门级天文望远镜、双筒望远镜及一般目视观察系统。

本次设计以有效焦距30 mm、入瞳直径4 mm、最大半视场角\(\pm 22.5^\circ\)为主要技术指标，采用F、d、C三条可见光谱线进行复色光分析。首先根据经典凯尔纳目镜结构建立初始模型，并设置0、0.25、0.50、0.75、1.00五个归一化视场，以保证轴上、中间视场和边缘视场均参与优化。随后在Merit Function中以RMS spot x+y centroid为核心评价指标，并加入有效焦距、系统尺度、光程和空气/玻璃厚度边界等约束，通过阻尼最小二乘法（DLS）局部优化、锤形优化（Hammer Optimization）全局搜索及人工微调，对透镜曲率半径、厚度、间隔和材料组合进行了系统优化。

最终设计结果表明，系统有效焦距为30.0 mm，后焦距约为20.950 mm，光阑面至像面的有限厚度合计约为48.303 mm，入瞳直径为4 mm，最大径向视场为22.5°。优化后系统在点列图、MTF、场曲畸变、相对照度、OPD波前误差和图像模拟等方面均表现良好：最大畸变约为1.7\%，边缘视场相对照度约为88\%，20 lp/mm处MTF约为0.71，轴上及全视场RMS波前误差均满足设计评价要求。公差分析结果表明，在所设定的曲率半径、中心厚度、折射率、阿贝数、偏心和倾斜公差范围内，系统仍具有较好的成像稳定性和工程实现可行性。

关键词：凯尔纳目镜；Zemax仿真；光学优化；像差分析；公差分析
